\section{Problem}

Queueing Theory is the study of waiting and cuts through in a fundamental way
across many disciplines. In our context we will use it to provide a foundation
for understanding performance in distributed systems. In some situations queues
will be explicit to the design our code, and in other it will be implicit to the
systems that we are building on.

For \textbf{one} (1) of the following resource sets I want you to watch,
summarize and reflect on the content. Additionally, I want you to synthesize at
least one questions. We will use these questions to start a discussion in class.
You will submit your reflections as a pdf file.

\begin{itemize}
    \item \textbf{Set 0}:
          \\
          \href
             {https://www.youtube.com/watch?v=rBIQmwaoZfs}
             {Queuing theory and Poisson process (25min, theoretical)}
          \\
          \href
             {https://www.youtube.com/watch?v=Hda5tMrLJqc}
             {%
                LISA17 - Queueing Theory in Practice: Performance Modeling for
                the Working Engineer (45min, practical)%
            }
            
    \item \textbf{Set 1}:
         \\
         \href
            {https://www.youtube.com/watch?v=rT7WR89GMC4}
            {%
                [Harchol-Balter] ``What Analytical Performance Modeling Teaches
                Us About Computer Systems Design.''%
            }
            \textit{(older video, but throrough)}
         
    \item \textbf{Set 2}:
          \\
          \href
            {https://www.youtube.com/watch?v=RXZRx_ET3v0}
            {%
                Mor Harchol-Balter: Optimal Scheduling for Multi-Server Systems
                (1hr)%
            }
    
    \item \textbf{Set 3}:
          \\
          \href
            {https://www.youtube.com/watch?v=lDn5vT6ymQs}
            {Fall 2020 LIDS Seminar — Mor Harchol-Balter (CMU) (1hr10min)}

    \item \textbf{Set 4}
          \\
          \href
            {https://www.youtube.com/watch?v=3at6ijBU2ug}
            {%
                SREcon24 Americas - System Performance and Queuing Theory -
                Concepts and Application (Intro-ish) (1hr)%
            }
\end{itemize}

\newpage
\section{Solution}

I chose the first set of videos.

After watching the theoretical video, I initially had difficulty connecting the
concepts to real-world engineering. Much of the video focused on formulas and
their derivations. Even though there were a few examples, the ideas did not
fully click for me at first.

However, after watching Eben Freeman's practical talk, \textit{Queueing Theory
in Practice: Performance Modeling for the Working Engineer}, I began to
understand why these concepts are useful. What surprised me the most was seeing
how the behavior of a system changes as more parallelism is introduced.

Throughout my experience, I had generally thought about horizontal scaling as
adding more instances in order to gain more processing capacity. The talk made
me realize that this relationship is not necessarily linear. Adding more workers
can initially improve throughput, but communication, synchronization, task
assignment, and other coordination costs can eventually reduce those benefits.
At a certain point, adding more resources may provide very little improvement or
can even make performance worse.

The load-balancing example made this particularly clear to me. If we try to find
the optimal worker for every task by checking every available worker, the cost
of making that decision grows along with the size of the system. A better
approach may be to use an approximation, such as randomly selecting two workers
and assigning the task to the better candidate. We give up the theoretically
optimal assignment in exchange for dramatically lower coordination overhead.

The distributed-query example also demonstrated an important architectural
principle. Dividing work among more nodes makes the individual scanning
operations faster, but collecting and aggregating the results introduces
additional work. Instead of having one node coordinate with every worker
directly, a hierarchy of intermediate aggregators can reduce this coordination
cost and allow the system to scale more efficiently.

My biggest realization from the video was that \textbf{horizontal scaling is not
automatically the solution to a performance problem}. Before adding more compute
resources, engineers should understand or attempt to understand where the
bottleneck actually exists. Sometimes profiling the application and reducing the
amount of work performed by each request can provide a greater improvement than
adding additional servers. In other cases, the architecture may need to change
so that less coordination is required between workers.

This also has a direct impact on infrastructure cost. If a performance problem
can be solved through code optimization, better task distribution, or a more
efficient architecture, an organization may avoid unnecessary horizontal or
vertical scaling.


My biggest takeaways are:
\begin{itemize}
    \item Think carefully about scalability while designing a system rather than
          assuming that adding resources will solve performance problems.
    
    \item Analyze the expected user load, throughput, latency, and computational
          requirements.

    \item Identify whether expensive tasks can be simplified or optimized before
          increasing compute capacity.
    
    \item Use small experiments, benchmarks, and simulations to estimate how a
          system will behave under greater load.
    
    \item Account for coordination overhead when parallelizing work.

    \item Consider approximate algorithms when finding a perfectly optimal
          solution would require excessive coordination.

    \item Measure real systems because mathematical models are useful only when
          their assumptions reasonably match reality.

    \item Ask engineering questions such as:
          \begin{itemize}
                \item \textit{Is this solution efficient in terms of both
                      latency and throughput?}
                \item \textit{What will happen if I double the number of
                      workers?}
                \item \textit{Can I reduce the amount of work instead of adding
                      more compute?}
          \end{itemize}
\end{itemize}

\section{Questions}
\textbf{Q}: What term describes a state in which increasing the workload or
concurrency causes a system to spend so many resources on coordination or
overhead that its useful throughput begins to decrease?
\\
\textbf{A}: Thrashing.

\textbf{Q}: What are some techniques that could help prevent or mitigate this
type of performance degradation?
\\
\textbf{A}: Reduce coordination, add throttling, etc.
